% Addendum to the existing service-message discussion; preserve section numbering.
\label{sec:security-boundaries}
\textbf{Authorization scope and credential reuse.} ABE restricts descriptor access to entities whose DKEY policies permit decryption, including authorized providers that are not selected. Names and traffic patterns remain visible. ACK and Response content encrypted only for a service permission can also be read by other holders of that permission.

ABE policies carry service permissions, allowing baseline trust rules to focus on identity and message-signing relationships rather than per-entity service-permission lists. Identity credentials and these rules can be reused across applications sharing the same identity trust configuration. Trust rules still specify which keys may sign each message type~\cite{yu2015trust}; when processing Selection, providers must also check that it comes from the requesting user, refers to the same service and request, contains the stored provider challenge, and is allowed by the current request state. Rejecting replay requires tracking whether the request has expired or the challenge has already been used. Issuing a replacement DKEY does not invalidate an old DKEY retained by its holder. Revocation requires key rotation or additional permission checks and cannot prevent use of previously disclosed plaintext.


Table~\ref{tab:message-acceptance} summarizes the checks required by the basic invocation flow. All four message types require signature validation, an authorized signing key for the message name, and consistent service and request identifiers. The challenge checks assume that authorized participants do not disclose protected values to unauthorized parties. The invocation experiments do not independently test every rejection condition.

\begin{table*}[t]
\centering
\caption{Acceptance checks for basic invocation messages}
\label{tab:message-acceptance}
\small
\begin{tabularx}{\textwidth}{@{}lX@{}}
\toprule
Message & Additional checks before advancing the invocation \\
\midrule
Request & The provider can decrypt content protected for the requested service, and evaluates whether to participate. A duplicate request must not create a conflicting challenge or execution. \\
ACK & The user matches RequestNonce to its active request and associates the feedback and ProviderChallenge with the signing provider. Only positive ACKs passing these checks supply candidates for selection. \\
Selection & The provider matches the requesting user, service, request, and stored ProviderChallenge. The request must still permit execution, and a repeated Selection must not start a second execution. \\
Response & The user matches RequestNonce, the active request, and a selected provider. Duplicate or late results must not be accepted as results of another request. \\
\bottomrule
\end{tabularx}
\end{table*}
